\EMhold

\EMsection{مقدمه}{مقدمه}

\begin{EMdefinition}{معادلهٔ دیفرانسیل}

معادله‌ای که در آن تابع و مشتقات وجود داشته باشند.

\end{EMdefinition}

\begin{EMdefinition}{معادلهٔ دیفرانسیل معمولی}

معادله‌ای که بر حسب فقط یک متغیر مستقل و مشتقات تابع فقط بر حسب آن متغیر نوشته شده باشد.

\end{EMdefinition}

\begin{EMdefinition}{معادلهٔ دیفرانسیل با مشتقات جزئی}

معادلهٔ دیفرانسیلی که تابع به چند متغیر مستقل وابسته بوده، مشتقات جزئی تابع بر حسب متغیرهای مختلف در معادله وجود داشته باشد.

\end{EMdefinition}

\begin{EMunit}

\textbf{معادلهٔ دیفرانسیل موج:}
\EMdisp{\frac{\partial^2u}{\partial t^2}\EMbr =c^2\frac{\partial^2u}{\partial x^2},\qquad\EMbr  \frac{\partial^2u}{\partial t^2}\EMbr =c^2\left(\frac{\partial^2u}{\partial x^2}+\frac{\partial^2u}{\partial y^2}\right),\qquad\EMbr  \frac{\partial^2u}{\partial t^2}\EMbr =c^2\left(\frac{\partial^2u}{\partial x^2}+\frac{\partial^2u}{\partial y^2}+\frac{\partial^2u}{\partial z^2}\right)}

\textbf{معادلهٔ دیفرانسیل انتقال حرارت:}
\EMdisp{\frac{\partial u}{\partial t}\EMbr =c^2\frac{\partial^2u}{\partial x^2},\qquad\EMbr  \frac{\partial u}{\partial t}\EMbr =c^2\left(\frac{\partial^2u}{\partial x^2}+\frac{\partial^2u}{\partial y^2}\right),\qquad\EMbr  \frac{\partial u}{\partial t}\EMbr =c^2\left(\frac{\partial^2u}{\partial x^2}+\frac{\partial^2u}{\partial y^2}+\frac{\partial^2u}{\partial z^2}\right)}

\end{EMunit}

\textbf{معادلهٔ دیفرانسیل لاپلاس:}
\EMdisp{\frac{\partial^2u}{\partial x^2}+\frac{\partial^2u}{\partial y^2}\EMbr =0,\qquad\EMbr  \frac{\partial^2u}{\partial x^2}+\frac{\partial^2u}{\partial y^2}+\frac{\partial^2u}{\partial z^2}\EMbr =0,\qquad\EMbr  \nabla^2u\EMbr =0}

\EMhold

\EMsection{معادلهٔ دیفرانسیل موجی یک‌بعدی}{معادلهٔ دیفرانسیل موجی یک‌بعدی}

\begin{EMunit}

مدلسازی یک تار مرتعش (به‌دست آوردن معادلهٔ حاکم بر \EMim{u(x,t)}).

\textbf{فرضیات فیزیکی:}

\begin{itemize}
\tightlist
\item
  جرم به صورت یکنواخت در طول تار توزیع شده (تار همگن)، کاملاً انعطاف‌پذیر است.
\item
  تار چنان محکم کشیده شده است که نیروی جاذبه در آن قابل چشم‌پوشی است.
\item
  میزان انحراف تار از حالت تعادل بسیار کوچک است.
\end{itemize}

\EMfigure{m04-string}{قطعهٔ \EMim{PQ} از تار مرتعش: کشش‌های \EMim{T_1} و \EMim{T_2} در دو انتها
با زوایای \EMim{\alpha} و \EMim{\beta}}

\end{EMunit}

\begin{EMunit}

\EMdisp{T_1\cos\alpha\EMbr =T_2\cos\beta\EMbr =T\EMbr =\text{const.}}
\EMdisp{T_2\sin\beta-T_1\sin\alpha\EMbr =\rho\,\Delta s\,\frac{\partial^2u}{\partial t^2}}
که در آن \EMim{\rho} چگالی جرمی تار و \EMim{\Delta s} دیفرانسیل طول تار است. با تقسیم بر \EMim{T}:
\EMdisp{T\big(\tan\beta-\tan\alpha\big)\EMbr =\rho\,\Delta s\,\frac{\partial^2u}{\partial t^2}\quad\EMbr \EMbr \Longrightarrow\quad\EMbr  T\left(\frac{\partial u}{\partial x}\bigg|_{(x+\Delta x,t)}-\frac{\partial u}{\partial x}\bigg|_{(x,t)}\right)\EMbr =\rho\,\Delta s\,\frac{\partial^2u}{\partial t^2}}

با استفاده از قضیهٔ مقدار میانی (\EMim{f(b)-f(a)=f'(l)(b-a)}، \EMim{l\in[a,b]}):
\EMdisp{T\left(\frac{\partial u}{\partial x}\bigg|_{(x+\Delta x,t)}-\frac{\partial u}{\partial x}\bigg|_{(x,t)}\right)\EMbr =T\frac{\partial^2u}{\partial x^2}\bigg|_{(l,t)}\Delta x\EMbr =\rho\,\Delta s\,\frac{\partial^2u}{\partial t^2},\qquad\EMbr  x<l<x+\Delta x}
طبق فرض \EMim{\displaystyle\lim_{\Delta x\to0}\frac{\Delta x}{\Delta s}=1}، پس:
\EMdisp{T\frac{\partial^2u}{\partial x^2}\EMbr =\rho\frac{\partial^2u}{\partial t^2},\qquad\EMbr  c^2\triangleq\frac T\rho\quad\EMbr \EMbr \Longrightarrow\quad\EMbr \frac{\partial^2u}{\partial t^2}\EMbr =c^2\frac{\partial^2u}{\partial x^2}}

\end{EMunit}

\begin{EMunit}

برای حل یکتای این معادله نیاز به شرایط مرزی و شرایط اولیه داریم:

\begin{itemize}
\tightlist
\item
  شرایط مرزی (\lr{Boundary Conditions}): \EMim{u(0,t)=u(L,t)=0}
\item
  شرایط اولیه (\lr{Initial Conditions}): \EMim{u(x,0)=f(x)} و \EMim{u_t(x,0)=g(x)}، که در آن \EMim{u_t\triangleq\dfrac{\partial u}{\partial t}}
\end{itemize}

\end{EMunit}

\EMhold

\EMsection{حل معادلهٔ دیفرانسیل موجی یک‌بعدی}{حل معادلهٔ دیفرانسیل موجی یک‌بعدی}

\begin{EMunit}

\EMdisp{\begin{cases}\dfrac{\partial^2u}{\partial t^2}=c^2\dfrac{\partial^2u}{\partial x^2}\\\noalign{\vskip 2mm} u(0,t)=u(L,t)=0\\ u(x,0)=f(x)\\ u_t(x,0)=g(x)\end{cases}}

\end{EMunit}

\begin{EMunit}

در اینجا برای حل این معادله از روش \textbf{جداسازی متغیرها} استفاده می‌کنیم. فرض:
\EMdisp{u(x,t)\EMbr =X(x)\times T(t)}
\EMdisp{\frac{\partial^2u}{\partial x^2}\EMbr =X''(x)T(t),\qquad\EMbr  \frac{\partial^2u}{\partial t^2}\EMbr =X(x)T''(t)\quad\EMbr \EMbr \Longrightarrow\quad\EMbr  X(x)T''(t)\EMbr =c^2X''(x)T(t)\quad\EMbr \EMbr \Longrightarrow\quad\EMbr \frac{T''}{c^2T}\EMbr =\frac{X''}{X}\EMbr =k}

\EMim{k} باید حتماً یک عدد منفی باشد. بررسی حالت‌ها:

\end{EMunit}

\begin{EMunit}

\textbf{اگر \EMim{k=0}:}
\EMdisp{\frac{X''}{X}\EMbr =0\ \EMbr \Rightarrow\ X''\EMbr =0\ \EMbr \Rightarrow\ X(x)\EMbr =Ax+B,\qquad\EMbr  X(0)\EMbr =0\EMbr \Rightarrow B\EMbr =0,\quad\EMbr  X(L)\EMbr =0\EMbr \Rightarrow A\EMbr =0\ \EMbr \Rightarrow\ X(x)\EMbr =0}

\textbf{اگر \EMim{k>0} (\EMim{k=\lambda^2}):}
\EMdisp{\frac{X''}{X}\EMbr =\lambda^2\ \EMbr \Rightarrow\ X''-\lambda^2X\EMbr =0\ \EMbr \Rightarrow\ X(x)\EMbr =Ae^{\lambda x}+Be^{-\lambda x},\qquad\EMbr  X(0)\EMbr =0\EMbr \Rightarrow A+B\EMbr =0,\quad\EMbr  X(L)\EMbr =0\EMbr \Rightarrow Ae^{\lambda L}+Be^{-\lambda L}\EMbr =0}
\EMdisp{\Longrightarrow\ A\EMbr =B\EMbr =0\ \EMbr \Rightarrow\ X(x)\EMbr =0}

\end{EMunit}

\begin{EMunit}

بنابراین \EMim{k} باید حتماً یک عدد منفی باشد (\EMim{k=-\lambda^2}):
\EMdisp{\frac{T''}{c^2T}\EMbr =\frac{X''}{X}\EMbr =-\lambda^2\quad\EMbr \EMbr \Longrightarrow\quad\EMbr \begin{cases}X''+\lambda^2X=0\\ T''+c^2\lambda^2T=0\end{cases}}
\EMdisp{X(x)\EMbr =A\cos\lambda x+B\sin\lambda x,\qquad\EMbr  X(0)\EMbr =0\EMbr \Rightarrow A\EMbr =0}
\EMdisp{X(L)\EMbr =0\ \EMbr \Rightarrow\ B\sin\lambda L\EMbr =0\ \EMbr \Rightarrow\ \lambda L\EMbr =n\pi\ \ (n\in\mathbb{Z})\ \EMbr \Rightarrow\ \lambda\EMbr =\frac{n\pi}L\EMbr =\lambda_n}
\EMdisp{X_n(x)\EMbr =B_n\sin\lambda_nx,\qquad\EMbr  \lambda_n\EMbr =\frac{n\pi}L\ \ (n\in\mathbb{Z})}
\EMdisp{T''+c^2\lambda_n^2T\EMbr =0\ \EMbr \Rightarrow\ T_n(t)\EMbr =C_n\cos(c\lambda_nt)+D_n\sin(c\lambda_nt)}
\EMdisp{u_n(x,t)\EMbr =X_n(x)\times T_n(t)\EMbr =B_n\sin\lambda_nx\times\big(C_n\cos(c\lambda_nt)+D_n\sin(c\lambda_nt)\big)}
با برهم‌نهی جواب‌ها:
\EMdisp{u(x,t)\EMbr =\sum_{n=1}^{\infty}\left[a_n\cos\!\left(\frac{cn\pi}Lt\right)+b_n\sin\!\left(\frac{cn\pi}Lt\right)\right]\sin\frac{n\pi}Lx}

\textbf{اعمال شرایط اولیه.} شرط \EMim{u(x,0)=f(x)}:
\EMdisp{u(x,0)\EMbr =\sum_{n=1}^{\infty}a_n\sin\frac{n\pi}Lx\EMbr =f(x)}
تابع \EMim{f(x)} با یک تابع تناوبی و فرد برابر شده است، بنابراین برای محاسبهٔ ضرایب مجهول فوریه باید \EMim{f(x)} را یک تابع تناوبی (با دورهٔ تناوب \EMim{2L}) و فرد فرض کنیم. به این صورت خواهیم داشت:
\EMdisp{a_n\EMbr =\frac2{2L}\int_{-L}^{L}f(x)\sin\!\left(\frac{2\pi}{2L}nx\right)dx\EMbr =\frac1L\int_{-L}^{L}f(x)\sin\!\left(\frac{n\pi}Lx\right)dx\EMbr =\frac2L\int_0^Lf(x)\sin\!\left(\frac{n\pi}Lx\right)dx}
شرط \EMim{u_t(x,0)=g(x)}:
\EMdisp{u_t(x,t)\EMbr =\sum_{n=1}^{\infty}\left[-a_n\frac{cn\pi}L\sin\!\left(\frac{cn\pi}Lt\right)+b_n\frac{cn\pi}L\cos\!\left(\frac{cn\pi}Lt\right)\right]\sin\frac{n\pi}Lx}
\EMdisp{u_t(x,0)\EMbr =\sum_{n=1}^{\infty}b_n\frac{cn\pi}L\sin\frac{n\pi}Lx\EMbr =g(x)\quad\EMbr \EMbr \Longrightarrow\quad\EMbr  b_n\EMbr =\frac2{cn\pi}\int_0^Lg(x)\sin\!\left(\frac{n\pi}Lx\right)dx}

\end{EMunit}

\begin{EMimportant}{جواب معادلهٔ موجی یک‌بعدی (روش جداسازی متغیرها)}

\EMdisp{u(x,t)\EMbr =\sum_{n=1}^{\infty}\left[a_n\cos\!\left(\frac{cn\pi}Lt\right)+b_n\sin\!\left(\frac{cn\pi}Lt\right)\right]\sin\frac{n\pi}Lx}
\EMdisp{a_n\EMbr =\frac2L\int_0^Lf(x)\sin\!\left(\frac{n\pi}Lx\right)dx,\qquad\EMbr  b_n\EMbr =\frac2{cn\pi}\int_0^Lg(x)\sin\!\left(\frac{n\pi}Lx\right)dx}

\end{EMimportant}

\begin{EMexample}{مثال ۱}

جواب معادلهٔ موجی یک‌بعدی زیر را به‌دست آورید.
\EMdisp{\frac{\partial^2u}{\partial t^2}\EMbr =\frac{\partial^2u}{\partial x^2},\qquad\EMbr  u(0,t)\EMbr =u(\pi,t)\EMbr =0,\qquad\EMbr  u(x,0)\EMbr =k\sin2x,\qquad\EMbr  u_t(x,0)\EMbr =0}

با توجه به اینکه سرعت اولیهٔ تار مرتعش صفر است، \EMim{g(x)=0\Rightarrow b_n=0}. بنابراین (\EMim{c=1}، \EMim{L=\pi}):
\EMdisp{u(x,t)\EMbr =\sum_{n=1}^{\infty}a_n\cos\!\left(\frac{cn\pi}Lt\right)\sin\frac{n\pi}Lx\EMbr =\sum_{n=1}^{\infty}a_n\cos nt\,\sin nx}
\EMdisp{a_n\EMbr =\frac2L\int_0^Lf(x)\sin\!\left(\frac{n\pi}Lx\right)dx\EMbr =\frac2\pi\int_0^\pi k\sin2x\sin nx\,dx\EMbr =\begin{cases}0&n\neq2\\ k&n=2\end{cases}}
\EMdisp{u(x,t)\EMbr =k\sin2x\cos2t}

\end{EMexample}

\begin{EMunit}

\EMfigure{m04-mode2}{مود دوم ارتعاش تار: \EMim{u(x,t)=k\sin2x\cos2t} در لحظه‌های مختلف}

\end{EMunit}

\begin{EMexample}{مثال ۲}

جواب معادلهٔ موجی یک‌بعدی زیر را به‌دست آورید.
\EMdisp{\frac{\partial^2u}{\partial t^2}\EMbr =\frac{\partial^2u}{\partial x^2},\qquad\EMbr  u(0,t)\EMbr =u(\pi,t)\EMbr =0,\qquad\EMbr  u(x,0)\EMbr =f(x),\qquad\EMbr  u_t(x,0)\EMbr =0}
که در آن \EMim{f(x)} تابع مثلثی شکل زیر است: \EMim{f(x)=\dfrac kax} برای \EMim{0\le x\le a} و \EMim{f(x)=\dfrac k{a-\pi}(x-\pi)} برای \EMim{a\le x\le\pi}.

\EMim{u_t(x,0)=g(x)=0\Rightarrow b_n=0} و
\EMdisp{u(x,t)\EMbr =\sum_{n=1}^{\infty}a_n\cos nt\,\sin nx}
\EMdisp{\begin{aligned}
a_n&=\frac2\pi\int_0^\pi f(x)\sin nx\,dx=\frac2\pi\int_0^a\frac ka\,x\sin nx\,dx+\frac2\pi\int_a^\pi\frac k{a-\pi}(x-\pi)\sin nx\,dx\\
&=\frac{2k}{a\pi}\left[x\left(\frac{-1}n\cos nx\right)-(1)\left(\frac{-1}{n^2}\sin nx\right)\right]_0^a+\frac{2k}{(a-\pi)\pi}\left[(x-\pi)\left(\frac{-1}n\cos nx\right)-(1)\left(\frac{-1}{n^2}\sin nx\right)\right]_a^\pi\\
&=\cdots=\frac{2k\sin na}{an^2(\pi-a)}
\end{aligned}}
\EMdisp{u(x,t)\EMbr =\frac{2k}{a(\pi-a)}\sum_{n=1}^{\infty}\frac{\sin na}{n^2}\cos nt\,\sin nx}

\end{EMexample}

\begin{EMunit}

\EMfigure{m04-triangle}{توزیع اولیهٔ مثلثی \EMim{f(x)} با قلهٔ \EMim{k} در \EMim{x=a}}

\end{EMunit}

\begin{EMexample}{مثال ۳}

جواب معادلهٔ موجی یک‌بعدی زیر را به‌دست آورید.
\EMdisp{\frac{\partial^2u}{\partial t^2}\EMbr =\frac{\partial^2u}{\partial x^2},\qquad\EMbr  u(0,t)\EMbr =u(\pi,t)\EMbr =0,\qquad\EMbr  u(x,0)\EMbr =kx(\pi-x),\qquad\EMbr  u_t(x,0)\EMbr =0}

\EMim{g(x)=0\Rightarrow b_n=0} و
\EMdisp{u(x,t)\EMbr =\sum_{n=1}^{\infty}a_n\cos nt\,\sin nx}
\EMdisp{a_n\EMbr =\frac2\pi\int_0^\pi kx(\pi-x)\sin nx\,dx\EMbr =\frac{2k}\pi\Big[(-x^2+\pi x)\Big(\frac{-1}n\cos nx\Big)-(-2x+\pi)\Big(\frac{-1}{n^2}\sin nx\Big)+(-2)\Big(\frac1{n^3}\cos nx\Big)\Big]_0^\pi\EMbr =\cdots\EMbr =\frac{4k\big(1-(-1)^n\big)}{\pi n^3}}
\EMdisp{u(x,t)\EMbr =\frac{4k}\pi\sum_{n=1}^{\infty}\frac{1-(-1)^n}{n^3}\cos nt\,\sin nx}

\end{EMexample}

\begin{EMunit}

\EMfigure{m04-parabola}{توزیع اولیهٔ سهمی‌شکل \EMim{f(x)=kx(\pi-x)}}

\end{EMunit}

\begin{EMexercise}{}

(الف) مقدار نسبت \EMim{\dfrac{a_1^2}{a_1^2+a_2^2+\cdots}} را به دست آورید.

(ب) نسبت \EMim{a_{2n-1}/a_{2n+1}} را بر حسب \EMim{n\in\mathbb{N}} رسم کنید.

\end{EMexercise}

\EMhold

\EMsection{روش جداسازی متغیرها: مثال‌ها}{روش جداسازی متغیرها: مثال‌ها}

\begin{EMexample}{مثال ۱}

مطلوب است حل معادلهٔ دیفرانسیل مقابل به روش جداسازی متغیرها: \EMim{xu_x=yu_y}.

\EMdisp{x\frac{\partial u}{\partial x}\EMbr =y\frac{\partial u}{\partial y},\qquad\EMbr  u(x,y)\EMbr =X(x)\cdot Y(y)}
\EMdisp{xX'Y\EMbr =yXY'\ \EMbr \Rightarrow\ \frac{xX'}X\EMbr =\frac{yY'}Y\EMbr =k}
\EMdisp{\frac{X'}X\EMbr =\frac kx\ \EMbr \Rightarrow\ \ln X\EMbr =k\ln x+c_1\ \EMbr \Rightarrow\ X(x)\EMbr =C_1x^k,\qquad\EMbr  \frac{Y'}Y\EMbr =\frac ky\ \EMbr \Rightarrow\ \ln Y\EMbr =k\ln y+c_2\ \EMbr \Rightarrow\ Y(y)\EMbr =C_2y^k}
\EMdisp{u(x,y)\EMbr =X(x)Y(y)\EMbr =C_1x^kC_2y^k\EMbr =C(xy)^k}

\end{EMexample}

\begin{EMexample}{مثال ۲}

مطلوب است حل معادلهٔ دیفرانسیل مقابل به روش جداسازی متغیرها: \EMim{u_{xx}+u_{yy}=0}.

\EMdisp{u(x,y)\EMbr =X(x)\cdot Y(y)\ \EMbr \Rightarrow\ X''Y+XY''\EMbr =0\ \EMbr \Rightarrow\ \frac{X''}X\EMbr =-\frac{Y''}Y\EMbr =k}
\EMdisp{\begin{cases}X''-kX=0\ \Rightarrow\ X(x)=Ae^{\sqrt kx}+Be^{-\sqrt kx}\\ Y''+kY=0\ \Rightarrow\ Y(y)=Ce^{\sqrt{-k}\,y}+De^{-\sqrt{-k}\,y}\end{cases}}
\EMdisp{u(x,y)\EMbr =X(x)Y(y)\EMbr =\big(Ae^{\sqrt kx}+Be^{-\sqrt kx}\big)\big(Ce^{\sqrt{-k}\,y}+De^{-\sqrt{-k}\,y}\big)}

\end{EMexample}

\begin{EMexample}{مثال ۳}

مطلوب است حل معادلهٔ دیفرانسیل مقابل به روش جداسازی متغیرها: \EMim{x^2u_{xy}=-3y^2u}.

\EMdisp{x^2\frac{\partial^2u}{\partial x\partial y}+3y^2u\EMbr =0,\qquad\EMbr  u(x,y)\EMbr =X(x)\cdot Y(y)\ \EMbr \Rightarrow\ x^2X'Y'+3y^2XY\EMbr =0\ \EMbr \Rightarrow\ \frac{x^2X'}X\EMbr =\frac{-3y^2Y}{Y'}\EMbr =k}
\EMdisp{\frac{X'}X\EMbr =\frac k{x^2}\ \EMbr \Rightarrow\ \ln X\EMbr =\frac{-k}x+c_1\ \EMbr \Rightarrow\ X(x)\EMbr =C_1e^{-k/x}}
\EMdisp{\frac{Y'}Y\EMbr =\frac{-3y^2}k\ \EMbr \Rightarrow\ \ln Y\EMbr =-\frac{y^3}k+c_2\ \EMbr \Rightarrow\ Y(y)\EMbr =C_2e^{-y^3/k}}
\EMdisp{u(x,y)\EMbr =X(x)Y(y)\EMbr =C_1e^{-k/x}\times C_2e^{-y^3/k}\EMbr =Ce^{-(k/x+y^3/k)}}

\end{EMexample}

\EMhold

\EMsection{معادلهٔ موجی یک‌بعدی با شرایط مرزی غیرصفر}{معادلهٔ موجی یک‌بعدی با شرایط مرزی غیرصفر}

\begin{EMexample}{مسئله}

جواب معادلهٔ دیفرانسیل موجی یک‌بعدی زیر را به‌دست آورید.
\EMdisp{\frac{\partial^2u}{\partial t^2}\EMbr =c^2\frac{\partial^2u}{\partial x^2},\qquad\EMbr  u(0,t)\EMbr =P_0,\quad\EMbr  u(L,t)\EMbr =P_1,\quad\EMbr  u(x,0)\EMbr =f(x),\quad\EMbr  u_t(x,0)\EMbr =g(x)}

\EMdisp{u(x,t)\EMbr =w(x,t)+v(x)\quad\EMbr \EMbr \Longrightarrow\quad\EMbr \begin{cases}w(0,t)+v(0)=P_0\\ w(L,t)+v(L)=P_1\end{cases}}
\EMdisp{\frac{\partial^2w}{\partial t^2}\EMbr =c^2\left(\frac{\partial^2w}{\partial x^2}+\frac{\partial^2v}{\partial x^2}\right)}
تابع \EMim{v(x)} را چنان انتخاب می‌کنیم که مسئلهٔ \EMim{w} شرایط مرزی همگن داشته باشد:
\EMdisp{\begin{cases}\dfrac{\partial^2v}{\partial x^2}=0\\ v(0)=P_0\\ v(L)=P_1\end{cases}\quad\EMbr \EMbr \Longrightarrow\quad\EMbr  v(x)\EMbr =\left(\frac{P_1-P_0}L\right)x+P_0}
\EMdisp{\begin{cases}\dfrac{\partial^2w}{\partial t^2}=c^2\dfrac{\partial^2w}{\partial x^2}\\ w(0,t)=w(L,t)=0\end{cases}\quad\EMbr \EMbr \Longrightarrow\quad\EMbr  w(x,t)\EMbr =\sum_{n=1}^{\infty}\left[A_n\cos\!\left(\frac{cn\pi}Lt\right)+B_n\sin\!\left(\frac{cn\pi}Lt\right)\right]\sin\frac{n\pi}Lx}
\EMdisp{u(x,t)\EMbr =\underbrace{\sum_{n=1}^{\infty}\left[A_n\cos\!\left(\frac{cn\pi}Lt\right)+B_n\sin\!\left(\frac{cn\pi}Lt\right)\right]\sin\frac{n\pi}Lx}_{w(x,t)}+\underbrace{\left(\frac{P_1-P_0}L\right)x+P_0}_{v(x)}}
شرط اولیهٔ \EMim{u(x,0)=f(x)}:
\EMdisp{\sum_{n=1}^{\infty}A_n\sin\frac{n\pi}Lx\EMbr =f(x)-\left(\frac{P_1-P_0}L\right)x-P_0\triangleq f_1(x)}
\EMdisp{A_n\EMbr =\frac2L\int_0^Lf_1(x)\sin\frac{n\pi}Lx\,dx\EMbr =\frac2L\int_0^L\left[f(x)-\left(\frac{P_1-P_0}L\right)x-P_0\right]\sin\frac{n\pi}Lx\,dx}

\end{EMexample}

\begin{EMexercise}{}

\EMim{B_n} را برای حالت (الف) \EMim{g(x)=0} و (ب) \EMim{g(x)} دلخواه به دست آورید.

\end{EMexercise}

\EMsection{معادلهٔ موجی یک‌بعدی با سرعت اولیهٔ صفر (امواج متحرک)}{معادلهٔ موجی یک‌بعدی با سرعت اولیهٔ صفر (امواج متحرک)}

اگر سرعت اولیه صفر باشد (\EMim{g(x)=0}) آنگاه خواهیم داشت:
\EMdisp{u(x,t)\EMbr =\sum_{n=1}^{\infty}a_n\cos\!\left(\frac{cn\pi}Lt\right)\sin\frac{n\pi}Lx,\qquad\EMbr  u(x,0)\EMbr =\sum_{n=1}^{\infty}a_n\sin\frac{n\pi}Lx\EMbr =f(x)\equiv\tilde f(x)}
با استفاده از \EMim{\sin\alpha\cos\beta=\frac12\big(\sin(\alpha+\beta)+\sin(\alpha-\beta)\big)}:
\EMdisp{\begin{aligned}
u(x,t)&=\sum_{n=1}^{\infty}a_n\frac12\left[\sin\!\left(\frac{n\pi}Lx+\frac{cn\pi}Lt\right)+\sin\!\left(\frac{n\pi}Lx-\frac{cn\pi}Lt\right)\right]\\
&=\frac12\sum_{n=1}^{\infty}a_n\left[\sin\frac{n\pi}L(x+ct)+\sin\frac{n\pi}L(x-ct)\right]\\
&=\frac12\sum_{n=1}^{\infty}a_n\sin\frac{n\pi}L(x+ct)+\frac12\sum_{n=1}^{\infty}a_n\sin\frac{n\pi}L(x-ct)=\frac12\Big[\tilde f(x+ct)+\tilde f(x-ct)\Big]
\end{aligned}}

\begin{EMimportant}{امواج متحرک}

\EMdisp{u(x,t)\EMbr =\frac12\Big[\tilde f(x+ct)+\tilde f(x-ct)\Big]}
\EMim{\tilde f} ادامهٔ فرد و متناوب (با دورهٔ \EMim{2L}) تابع اولیهٔ \EMim{f} است؛ جواب برهم‌نهی دو موج با نیمی از دامنهٔ اولیه است که با سرعت \EMim{c} در دو جهت مخالف حرکت می‌کنند.

\end{EMimportant}

\begin{EMunit}

\EMfigure{m04-traveling}{ادامهٔ فرد و متناوب \EMim{\tilde f(x)} و دو موج
\EMim{\frac12\tilde f(x+ct)} و \EMim{\frac12\tilde f(x-ct)} که در دو جهت
مخالف حرکت می‌کنند}

\end{EMunit}

\begin{EMexample}{}

جواب معادلهٔ موجی یک‌بعدی زیر را با فرض سرعت اولیهٔ صفر در زمان‌های مختلف رسم کنید.
\EMdisp{\frac{\partial^2u}{\partial t^2}\EMbr =c^2\frac{\partial^2u}{\partial x^2},\qquad\EMbr  u(0,t)\EMbr =u(L,t)\EMbr =0,\qquad\EMbr  u(x,0)\EMbr =f(x),\qquad\EMbr  u_t(x,0)\EMbr =0}
\EMdisp{f(x)\EMbr =\begin{cases}\dfrac{2k}Lx&0<x<\dfrac L2\\\noalign{\vskip 7mm} \dfrac{2k}L(L-x)&\dfrac L2<x<L\end{cases}}

با رسم \EMim{\frac12\tilde f(x+ct)} و \EMim{\frac12\tilde f(x-ct)} و جمع آنها در زمان‌های \EMim{t=0}، \EMim{t=L/5c}، \EMim{t=2L/5c}، \EMim{t=L/2c} و \EMim{t=3L/5c} شکل جواب به دست می‌آید.

\end{EMexample}

\begin{EMunit}

\EMfigure{m04-triangle-time}{جواب مثال برای \EMim{t=0}، \EMim{L/5c}، \EMim{2L/5c}، \EMim{L/2c} و
\EMim{3L/5c}؛ در \EMim{t=L/2c} شکل تار صاف است}

\end{EMunit}
